% Original analytic teaching diagram, 2026-11-12.
% Recording two dice without order merges the 36 ordered outcomes into 21 classes.
% Diagonal (doubles, blue): each class has 1 ordered outcome, mass 1/36.
% Off-diagonal classes have 2 mirror cells, mass 2/36; examples {1,2} (yellow) and {3,5} (purple).
\definecolor{VB}{HTML}{7188A3}\definecolor{VBF}{HTML}{E5EAF0}
\definecolor{VG}{HTML}{829C8B}\definecolor{VGF}{HTML}{E8EEE9}
\definecolor{VY}{HTML}{B4A16F}\definecolor{VYF}{HTML}{F1ECDD}
\definecolor{VR}{HTML}{A7787D}\definecolor{VRF}{HTML}{F0E5E6}
\definecolor{VP}{HTML}{9588A5}\definecolor{VPF}{HTML}{ECE8F0}
\definecolor{VI}{HTML}{3B4252}\definecolor{VM}{HTML}{6B7280}\definecolor{VL}{HTML}{4C4D4F}
\begin{tikzpicture}[x=1mm,y=1mm,>=Stealth,font=\figurefont\fontsize{8}{10}\selectfont,
lab/.style={text=FigureMuted,font=\figurefont\fontsize{7.5}{9}\selectfont},
txt/.style={text=VI},
cell/.style={draw=VL,line width=.5pt},
ev/.style={draw=VB,fill=VBF,line width=.7pt},
yc/.style={draw=VY,fill=VYF,line width=.9pt},
pc/.style={draw=VP,fill=VPF,line width=.9pt}]
\def\s{6}
\foreach \i in {1,...,6}{\foreach \j in {1,...,6}{\draw[cell] ({(\i-1)*\s},{(\j-1)*\s}) rectangle ++(\s,\s);}}
\foreach \k in {1,...,6}{\draw[ev] ({(\k-1)*\s},{(\k-1)*\s}) rectangle ++(\s,\s);}
\draw[yc] (0,\s) rectangle ++(\s,\s); \draw[yc] (\s,0) rectangle ++(\s,\s);
\draw[pc] ({2*\s},{4*\s}) rectangle ++(\s,\s); \draw[pc] ({4*\s},{2*\s}) rectangle ++(\s,\s);
\draw[VM,line width=.5pt,dash pattern=on 2pt off 1.5pt,<->] ({1.5*\s},{0.5*\s}) to[bend right=40] ({0.5*\s},{1.5*\s});
\draw[VM,line width=.5pt,dash pattern=on 2pt off 1.5pt,<->] ({4.5*\s},{2.5*\s}) to[bend right=30] ({2.5*\s},{4.5*\s});
\foreach \k in {1,...,6}{\node[lab] at ({(\k-.5)*\s},-2.8) {\k}; \node[lab] at (-2.8,{(\k-.5)*\s}) {\k};}
\node[lab] at ({3*\s},-7) {第一枚};
\node[lab,rotate=90] at (-7,{3*\s}) {第二枚};
\end{tikzpicture}
